% ============================================================================
%  Laboratorio di Programmazione -- Lezione 12 (24 novembre 2026)
%  Compilare con:  pdflatex -shell-escape Lezione12_Comparable_Collezioni_ADT.tex
%  Richiede: mermaid.sty (nel progetto) + Mermaid CLI `mmdc` nel PATH
%  NB: ogni frame che contiene codice (lstlisting) o un diagramma (mermaid)
%      deve essere dichiarato [fragile].
% ============================================================================
\documentclass[10pt,aspectratio=169,xcolor={table}]{beamer}
\usetheme{Madrid}
\usecolortheme{default}

% `provide=*` carica la localizzazione italiana dai file .ini di babel: funziona anche
% su installazioni TeX prive di texlive-lang-italian; con una TeX Live completa e`
% equivalente al classico \usepackage[italian]{babel}.
\usepackage[italian,provide=*]{babel}
\usepackage[utf8]{inputenc}

% Localizzazione dei nomi dei blocchi Beamer
\uselanguage{italian}
\languagepath{italian}
\deftranslation[to=italian]{Definition}{Definizione}
\deftranslation[to=italian]{Theorem}{Teorema}
\deftranslation[to=italian]{Example}{Esempio}
\deftranslation[to=italian]{Proof}{Dimostrazione}
\deftranslation[to=italian]{Corollary}{Corollario}
\deftranslation[to=italian]{Lemma}{Lemma}

\usepackage{listings}
\usepackage{xcolor}
\usepackage{hyperref}
\usepackage{array}
\usepackage{booktabs}
\usepackage{mermaid}

% Mermaid: sfondo trasparente e dimensione massima adatta a una slide 16:9
\MermaidRendererOptions{-b transparent}
\MermaidAdjustBoxOpts{max width=0.95\linewidth,max height=0.62\textheight,center}

% Configurazione colori per il codice
\definecolor{codegreen}{rgb}{0,0.6,0}
\definecolor{codegray}{rgb}{0.5,0.5,0.5}
\definecolor{codepurple}{rgb}{0.58,0,0.82}
\definecolor{backcolour}{rgb}{0.95,0.95,0.92}
\definecolor{keywordblue}{rgb}{0.0,0.0,0.7}

% Stile per snippet Java
\lstdefinestyle{javastyle}{
    backgroundcolor=\color{backcolour},
    commentstyle=\color{codegreen}\itshape,
    keywordstyle=\color{keywordblue}\bfseries,
    numberstyle=\scriptsize\color{codegray},
    stringstyle=\color{codepurple},
    basicstyle=\ttfamily\footnotesize,
    breakatwhitespace=false,
    breaklines=true,
    captionpos=b,
    keepspaces=true,
    numbers=none,
    numbersep=5pt,
    showspaces=false,
    showstringspaces=false,
    showtabs=false,
    tabsize=2,
    aboveskip=0.6\medskipamount,
    belowskip=0.6\medskipamount,
    language=Java,
    morekeywords={default},
    literate={à}{{\`a}}1 {è}{{\`e}}1 {é}{{\'e}}1 {ì}{{\`i}}1
             {ò}{{\`o}}1 {ù}{{\`u}}1 {È}{{\`E}}1
}
\lstset{style=javastyle}

% Stile per righe di terminale
\lstdefinestyle{shell}{
    language=bash,
    backgroundcolor=\color{black!85},
    basicstyle=\ttfamily\footnotesize\color{white},
    keywordstyle=\color{white},
    stringstyle=\color{white},
    commentstyle=\color{green!70},
    numbers=none,
    frame=none,
    aboveskip=0.6\medskipamount,
    belowskip=0.6\medskipamount,
    showstringspaces=false
}

% Celle dei diagrammi di pila e coda realizzate con array
\newcommand{\vuota}{\cellcolor{gray!15}}
\newcommand{\piena}[1]{\cellcolor{blue!12}\texttt{#1}}

\title[Lezione 12 -- Comparable, collezioni e ADT]{Lezione 12 -- \texttt{Comparable}, collezioni e ADT}
\subtitle{l'ordine naturale e i comparatori, l'ordinamento con \texttt{Arrays.sort}; pila e coda con un array, il nodo e la lista concatenata, pila e coda con una lista concatenata}
\author[Mattia Fonisto]{Mattia Fonisto}
\institute[Lab.\ di Programmazione]{Laboratorio di Programmazione -- Ingegneria Informatica, II anno\\
           Università degli Studi di Napoli Federico II -- Sede di San Giovanni a Teduccio}
\date{24 novembre 2026}

\begin{document}

\frame{\titlepage}

% ========================================================================
\begin{frame}{Agenda della lezione (2h)}
\begin{enumerate}
    \item \textbf{Ripasso e Homework 11} \hfill (\emph{15 min})
    \begin{itemize}
        \item uno di voi alla lavagna con \texttt{Intervallo}, \texttt{Occorrenze} o \texttt{Classifica}: lo commentiamo insieme
    \end{itemize}
    \item \textbf{\texttt{Comparable} e \texttt{Comparator}} \hfill (\emph{15 min})
    \begin{itemize}
        \item l'ordine naturale di una classe nostra; \texttt{Arrays.sort} su un array di oggetti
        \item un ordine diverso: i comparatori
    \end{itemize}
    \item \textbf{Collezioni e tipi di dato astratti} \hfill (\emph{30 min})
    \begin{itemize}
        \item contenitori e politiche di accesso; precondizioni e postcondizioni
        \item la pila e la coda con un array; il nodo e la lista concatenata; la pila con una lista concatenata
    \end{itemize}
    \item \textbf{Esercitazione} -- tutti insieme \hfill (\emph{50 min})
    \begin{itemize}
        \item \emph{Commentiamo insieme}: le parentesi bilanciate, con una pila
        \item \emph{Insieme, passo passo}: la coda con una lista concatenata
    \end{itemize}
    \item \textbf{Wrap-up e Homework 12} \hfill (\emph{10 min})
\end{enumerate}
\end{frame}

% ========================================================================
\section{Comparable e Comparator}

\begin{frame}[fragile]{\texttt{Comparable} sulle nostre classi: l'ordine naturale -- \texttt{Persona.java}}
\begin{columns}[T]
\column{0.5\textwidth}
\begin{lstlisting}[basicstyle=\ttfamily\scriptsize]
public class Persona
        implements Comparable<Persona> {

    private final String nome;
    private final String cognome;
    private final int eta;
\end{lstlisting}
\begin{lstlisting}[basicstyle=\ttfamily\scriptsize]
/** Negativo se this e' piu' giovane. */
@Override
public int compareTo(Persona altra) {
    return Integer.compare(eta, altra.eta);
}

@Override
public String toString() {
    return cognome + " " + nome
            + " (" + eta + ")";
}
\end{lstlisting}
\begin{lstlisting}[basicstyle=\ttfamily\scriptsize]
Persona p1 = new Persona("Anna", "Russo", 35);
Persona p2 = new Persona("Luca", "Verde", 20);
System.out.println(p1.compareTo(p2));   // 1
System.out.println(p2.compareTo(p1));   // -1
\end{lstlisting}

\column{0.5\textwidth}
\begin{block}{Ordine naturale}
\small Una classe che implementa \texttt{Comparable<T>} con \texttt{T} uguale a se stessa dà ai suoi oggetti un \emph{ordine naturale}: \texttt{a.compareTo(b)} è negativo se \texttt{a} viene prima di \texttt{b}, zero se stanno nello stesso posto, positivo se viene dopo.
\end{block}
\footnotesize
\begin{itemize}\itemsep=3pt
    \item \texttt{Comparable<T>} è l'interfaccia generica della lezione~11: con l'argomento \texttt{Persona}, il metodo da scrivere è \texttt{compareTo(Persona altra)}, e dentro si legge \texttt{altra.eta} (stessa classe, lezione~5)
    \item \textbf{l'ordine è una scelta}: qui per età. \texttt{Integer.compare} dà negativo, zero o positivo, come \texttt{Double.compare} (lezione~8); \texttt{eta - altra.eta} avrebbe lo stesso segno
    \item due coetanei danno \texttt{0} anche se \texttt{equals} direbbe \texttt{false}: conta solo ciò che ordina
    \item costruttore e \texttt{get} non mostrati; è una nuova \texttt{Persona}, con l'età
\end{itemize}
\end{columns}
\end{frame}

% ========================================================================
\begin{frame}[fragile]{Ordinare un array di oggetti: \texttt{Arrays.sort} -- \texttt{ProvaPersone.java}}
\begin{columns}[T]
\column{0.5\textwidth}
\begin{lstlisting}[basicstyle=\ttfamily\scriptsize]
Persona[] elenco = { p1, p2,
    new Persona("Mario", "Annunziata", 33),
    new Persona("Sara", "Esposito", 33) };
Arrays.sort(elenco);      // ordine naturale
for (Persona p : elenco) {
    System.out.println(p);
}
System.out.println(Utili.max(elenco)); // max
\end{lstlisting}
\begin{lstlisting}[style=shell,basicstyle=\ttfamily\scriptsize\color{white}]
Verde Luca (20)
Annunziata Mario (33)
Esposito Sara (33)
Russo Anna (35)
Russo Anna (35)
\end{lstlisting}
\footnotesize
\begin{itemize}\itemsep=3pt
    \item \texttt{Arrays.sort} ordina \textbf{sul posto} (lezione~4) chiamando il \texttt{compareTo} degli elementi: funziona su un array di qualunque classe che implementa \texttt{Comparable}; i due coetanei restano nell'ordine in cui erano
    \item il \texttt{max} generico della lezione~11, con il vincolo \texttt{<T extends Comparable<T>>}, ora accetta un \texttt{Persona[]}, come ogni metodo con quel vincolo: la promessa della scorsa lezione
\end{itemize}

\column{0.5\textwidth}
\footnotesize
\textbf{E se \texttt{Persona} non implementasse \texttt{Comparable}?} (\texttt{ProvaSenzaComparable.java})
\begin{lstlisting}[basicstyle=\ttfamily\scriptsize]
Arrays.sort(elenco);   // compila, ma...
\end{lstlisting}
\begin{lstlisting}[style=shell,basicstyle=\ttfamily\scriptsize\color{white}]
Exception in thread "main"
java.lang.ClassCastException: class Persona
cannot be cast to class java.lang.Comparable
(...)
\end{lstlisting}
\begin{itemize}\itemsep=3pt
    \item con \texttt{Utili.max(elenco)} l'errore uscirebbe a tempo di \textbf{compilazione}: il vincolo rifiuta un \texttt{T} non confrontabile. \texttt{Arrays.sort} si comporta diversamente perché la sua intestazione è \texttt{sort(Object[] a)}, senza vincolo: accetta qualunque array di oggetti e fa il cast a \texttt{Comparable} elemento per elemento \textbf{a run time}
    \item la regola: per ordinare con \texttt{Arrays.sort} o cercare con \texttt{Arrays.binarySearch} (lezione~4) un array di oggetti nostri, la classe deve implementare \texttt{Comparable}
\end{itemize}
\end{columns}
\end{frame}

% ========================================================================
\begin{frame}[fragile]{Un altro ordine: \texttt{Comparator} -- \texttt{ComparaPerCognome.java}}
\begin{columns}[T]
\column{0.5\textwidth}
\begin{lstlisting}[basicstyle=\ttfamily\scriptsize]
public interface Comparator<T> {
    int compare(T a, T b);   // come compareTo
}
\end{lstlisting}
\begin{lstlisting}[basicstyle=\ttfamily\scriptsize]
import java.util.Comparator;

public class ComparaPerCognome
        implements Comparator<Persona> {

    @Override
    public int compare(Persona a, Persona b) {
        String ca = a.getCognome();
        String cb = b.getCognome();
        return ca.compareTo(cb);
    }
}
\end{lstlisting}
\begin{block}{Comparatore}
\small Un \emph{comparatore} è un oggetto di una classe che implementa \texttt{Comparator<T>} (\texttt{java.util}, qui semplificata): di solito non ha dati, ha solo \texttt{compare(a, b)}, con le regole di segno di \texttt{compareTo}.
\end{block}

\column{0.5\textwidth}
\begin{lstlisting}[basicstyle=\ttfamily\scriptsize]
Arrays.sort(elenco, new ComparaPerCognome());
for (Persona p : elenco) {
    System.out.println(p);
}
ComparaPerCognome c = new ComparaPerCognome();
System.out.println(c.compare(p1, p2));
\end{lstlisting}
\begin{lstlisting}[style=shell,basicstyle=\ttfamily\scriptsize\color{white}]
Annunziata Mario (33)
Esposito Sara (33)
Russo Anna (35)
Verde Luca (20)
-4
\end{lstlisting}
\footnotesize
\begin{itemize}\itemsep=3pt
    \item l'ordine naturale è \textbf{uno solo}; per un altro ordine si passa ad \texttt{Arrays.sort} un secondo argomento, il comparatore: \texttt{sort(T[] v, Comparator<T> c)} (semplificata). Ogni ordine in più è una classe in più (\texttt{ComparaPerNome}, con \texttt{getNome})
    \item \texttt{compare} confronta i due \textbf{argomenti}, non \texttt{this}: il comparatore sta fuori da \texttt{Persona}, che non cambia, e usa i \texttt{get}; \texttt{-4} è \texttt{'R'} $-$ \texttt{'V'} (lezione~2)
\end{itemize}
\end{columns}
\end{frame}

% ========================================================================
\section{Collezioni e tipi di dato astratti}

\begin{frame}[fragile]{Collezioni: contenitori e politiche di accesso}
\begin{columns}[T]
\column{0.52\textwidth}
\begin{block}{Collezione}
\small Una \emph{collezione} è un oggetto che contiene altri oggetti, i suoi elementi; la classe che la realizza si chiama anche \emph{contenitore}. Le collezioni si distinguono per la \emph{politica di accesso}: in quale posizione va un elemento quando viene inserito, e quale elemento viene letto o tolto per primo.
\end{block}
\scriptsize
\begin{center}
\renewcommand{\arraystretch}{1.15}
\begin{tabular}{|>{\raggedright\arraybackslash}p{2.5cm}|>{\raggedright\arraybackslash}p{4.5cm}|}
\hline
\textbf{Politica} & \textbf{Chi decide la posizione di un elemento} \\
\hline
nessuna & il \textbf{client}: l'elemento va dove indica lui, come \texttt{v[3] = x} in un array \\
\hline
ordinata & il \textbf{valore} dell'elemento: 48 inserito in \texttt{[52, 55, 61]} va prima di 52, per mantenere l'ordine \\
\hline
LIFO (\emph{last in, first out}): la \textbf{pila} (\emph{stack}) & l'\textbf{ordine di arrivo}: esce per primo l'ultimo entrato (i piatti impilati) \\
\hline
FIFO (\emph{first in, first out}): la \textbf{coda} (\emph{queue}) & l'\textbf{ordine di arrivo}: esce per primo il primo entrato (la fila allo sportello) \\
\hline
\end{tabular}
\end{center}

\column{0.03\textwidth}
\column{0.45\textwidth}
\begin{block}{Precondizione e postcondizione}
\small La \emph{precondizione} di un metodo è ciò che deve essere vero quando viene chiamato; la \emph{postcondizione} è ciò che è vero quando termina: il risultato e l'effetto sullo stato. Insieme dicono che cosa fa il metodo, senza dire come.
\end{block}
\footnotesize
\begin{itemize}\itemsep=3pt
    \item precondizione violata: il metodo non promette nulla; le classi di oggi sollevano \texttt{IllegalStateException} (non controllata, come \texttt{IllegalArgumentException}): l'oggetto non è nello stato giusto. Il client la \textbf{evita} chiedendo prima, e non la cattura (lezione~10)
    \item pila e coda sono \textbf{tipi di dato astratti} (lezione~5): contano le operazioni, con pre e postcondizioni, non la struttura dati sotto
\end{itemize}
\end{columns}
\end{frame}

% ========================================================================
\begin{frame}[fragile]{L'ADT pila (\emph{stack}): LIFO e l'interfaccia \texttt{Pila<T>} -- \texttt{Pila.java}}
\begin{columns}[T]
\column{0.5\textwidth}
\footnotesize
\begin{center}
\renewcommand{\arraystretch}{1.15}
\begin{tabular}{|l|l|>{\raggedright\arraybackslash}p{3.6cm}|}
\hline
\textbf{Operazione} & \textbf{Pre} & \textbf{Post} \\
\hline
\texttt{push(x)} & non piena & \texttt{x} è in cima; un elemento in più \\
\hline
\texttt{pop()} & non vuota & restituisce l'elemento in cima e lo toglie \\
\hline
\texttt{isEmpty()} & -- & vero se non ci sono elementi \\
\hline
\texttt{isFull()} & -- & vero se non c'è posto per un altro \\
\hline
\end{tabular}
\end{center}
\begin{itemize}\itemsep=3pt
    \item \textbf{LIFO}: si inserisce e si toglie sempre dallo stesso lato, la \textbf{cima}; l'elemento che esce è quello entrato per \textbf{ultimo}, mentre il primo entrato esce per ultimo, come l'ultimo piatto messo sulla pila è il primo che si prende
    \item \textbf{l'interfaccia dell'ADT è un'\texttt{interface} di Java} (lezione~9): i metodi con le pre e postcondizioni nei commenti; il \textbf{corpo} è una classe che la implementa, e possono essercene più d'una. I nomi delle operazioni sono quelli inglesi, usati ovunque
\end{itemize}

\column{0.5\textwidth}
\begin{lstlisting}[basicstyle=\ttfamily\scriptsize]
/** L'ADT pila (stack): politica LIFO. */
public interface Pila<T> {

    /** Mette x in cima. Pre: !isFull(). */
    void push(T x);

    /** Toglie e restituisce la cima.
        Pre: !isEmpty(). */
    T pop();

    boolean isEmpty();

    boolean isFull();
}
\end{lstlisting}
\begin{center}
\scriptsize\renewcommand{\arraystretch}{1.2}\setlength{\tabcolsep}{4pt}
\begin{tabular}{|>{\raggedright\arraybackslash}p{2.3cm}|>{\raggedright\arraybackslash}p{1.6cm}|>{\raggedright\arraybackslash}p{2.3cm}|}
\hline
\textbf{operazione} & \textbf{esce} & \textbf{pila (cima a destra)} \\
\hline
\texttt{push} di Anna, Luca, Sara & & \texttt{Anna Luca Sara} \\
\hline
\texttt{pop()} & \texttt{Sara} & \texttt{Anna Luca} \\
\hline
\texttt{push(Marco)} & & \texttt{Anna Luca Marco} \\
\hline
tre \texttt{pop()} & \texttt{Marco}, \texttt{Luca}, \texttt{Anna} & vuota \\
\hline
\end{tabular}
\end{center}
\end{columns}
\end{frame}

% ========================================================================
\begin{frame}[fragile]{La pila con un array: lo stato -- \texttt{PilaArray.java}}
\begin{columns}[T]
\column{0.44\textwidth}
\begin{lstlisting}[basicstyle=\ttfamily\scriptsize]
public class PilaArray<T>
        implements Pila<T> {

    // non T[] (lezione 11)
    private final Object[] elementi;
    // prima posizione libera
    private int cima;

    public PilaArray(int capienza) {
        elementi = new Object[capienza];
        cima = 0;
    }

    @Override
    public boolean isEmpty() {
        return cima == 0;
    }

    @Override
    public boolean isFull() {
        return cima == elementi.length;
    }
\end{lstlisting}

\column{0.56\textwidth}
\footnotesize
\textbf{Dopo \texttt{push("Anna")}, \texttt{push("Luca")}, \texttt{push("Sara")} su una pila di capienza 5:}
\begin{center}
\renewcommand{\arraystretch}{1.3}\setlength{\tabcolsep}{5pt}
\begin{tabular}{l|c|c|c|c|c|}
\cline{2-6}
indice & 0 & 1 & 2 & 3 & 4 \\
\cline{2-6}
\texttt{elementi} & \piena{Anna} & \piena{Luca} & \piena{Sara} & \vuota & \vuota \\
\cline{2-6}
\multicolumn{1}{l}{} & \multicolumn{3}{l}{} & \multicolumn{2}{l}{$\uparrow$ \texttt{cima = 3}} \\
\end{tabular}
\end{center}
\begin{itemize}\itemsep=3pt
    \item \textbf{lo stato}: un array di \texttt{Object} (il rimedio della lezione~11 a \texttt{new T[n]}) di capienza fissa, e l'indice \texttt{cima}. La regola che vale sempre: gli elementi occupano le posizioni da \texttt{0} a \texttt{cima - 1}, la \textbf{cima della pila} è in \texttt{cima - 1} e \texttt{cima} è la prima posizione libera
    \item \texttt{isEmpty} e \texttt{isFull} non fanno che leggere questa regola: vuota se \texttt{cima} è 0, piena se \texttt{cima} ha raggiunto la lunghezza dell'array
    \item l'array è \texttt{final}: il riferimento non cambia mai, il contenuto sì (lezione~6); la capienza la sceglie il client nel costruttore
\end{itemize}
\end{columns}
\end{frame}

% ========================================================================
\begin{frame}[fragile]{La pila con un array: \texttt{push} e \texttt{pop} -- \texttt{PilaArray.java}}
\begin{columns}[T]
\column{0.54\textwidth}
\begin{lstlisting}[basicstyle=\ttfamily\scriptsize]
@Override
public void push(T x) {
    if (isFull()) {
        throw new IllegalStateException("piena");
    }
    elementi[cima] = x;
    cima++;
}

@Override
public T pop() {
    if (isEmpty()) {
        throw new IllegalStateException("vuota");
    }
    cima--;
    return (T) elementi[cima];
}
\end{lstlisting}
\footnotesize
\begin{itemize}\itemsep=3pt
    \item ogni metodo controlla la \textbf{sua} precondizione, poi mantiene la regola: \texttt{push} scrive nella prima posizione libera e la fa avanzare; \texttt{pop} fa il contrario, nell'ordine contrario: per questo esce l'ultimo entrato
    \item il cast \texttt{(T)} e l'avviso \emph{unchecked}: come nella lezione~11, la garanzia la dà \texttt{push}
\end{itemize}

\column{0.46\textwidth}
\begin{lstlisting}[basicstyle=\ttfamily\scriptsize]
Pila<String> p = new PilaArray<>(2);
p.push("Anna");
p.push("Luca");
System.out.println(p.isFull());
System.out.println(p.pop());
System.out.println(p.pop());
System.out.println(p.isEmpty());
\end{lstlisting}
\begin{lstlisting}[style=shell,basicstyle=\ttfamily\scriptsize\color{white}]
true
Luca
Anna
true
\end{lstlisting}
\footnotesize
\begin{itemize}\itemsep=3pt
    \item in \texttt{ProvaPile}, \texttt{p} è dichiarata di tipo \texttt{Pila<String>}, l'\textbf{interfaccia}, non \texttt{PilaArray<String>}: così il codice che usa \texttt{p} può chiamare solo le operazioni dell'ADT e non dipende dalla realizzazione; per passare a un'altra basta cambiare il \texttt{new}, e il resto non si tocca (lezione~9). Il \texttt{new}, invece, deve nominare una classe: un'interfaccia non si istanzia
    \item su una pila di 2 un terzo \texttt{push} solleverebbe l'eccezione
\end{itemize}
\end{columns}
\end{frame}

% ========================================================================
\begin{frame}[fragile]{L'ADT coda (\emph{queue}): FIFO e l'idea della coda circolare -- \texttt{Coda.java}}
\begin{columns}[T]
\column{0.5\textwidth}
\begin{lstlisting}[basicstyle=\ttfamily\scriptsize]
/** L'ADT coda (queue): politica FIFO. */
public interface Coda<T> {

    /** Mette x in fondo. Pre: !isFull(). */
    void enqueue(T x);

    /** Toglie e restituisce il primo.
        Pre: !isEmpty(). */
    T dequeue();

    boolean isEmpty();

    boolean isFull();
}
\end{lstlisting}
\footnotesize
\begin{itemize}\itemsep=3pt
    \item le stesse precondizioni della pila; cambiano le postcondizioni: si inserisce da una parte, \textbf{in fondo}, e si toglie dall'altra, \textbf{in testa}: il primo arrivato è il primo servito
    \item con un array servono \textbf{due indici}: \texttt{testa}, dove sta il primo, e \texttt{coda}, dove andrà il prossimo \texttt{enqueue}; \texttt{dequeue} fa avanzare \texttt{testa}, \texttt{enqueue} fa avanzare \texttt{coda}: i due indici si inseguono verso destra
\end{itemize}

\column{0.5\textwidth}
\footnotesize
\textbf{Coda circolare, capienza 4}: tre \texttt{enqueue}, \texttt{dequeue()}, due \texttt{enqueue}
\begin{center}
\scriptsize\renewcommand{\arraystretch}{1.3}\setlength{\tabcolsep}{3pt}
\begin{tabular}{l|c|c|c|c|l}
\cline{2-5}
indice & 0 & 1 & 2 & 3 & \\
\cline{2-5}
3 \texttt{enqueue} & \piena{Anna} & \piena{Luca} & \piena{Sara} & \vuota & \texttt{testa 0, coda 3} \\
\cline{2-5}
\texttt{dequeue()} & \vuota & \piena{Luca} & \piena{Sara} & \vuota & \texttt{testa 1, coda 3} \\
\cline{2-5}
\texttt{enqueue(Marco)} & \vuota & \piena{Luca} & \piena{Sara} & \piena{Marco} & \texttt{testa 1, coda 0} \\
\cline{2-5}
\texttt{enqueue(Elena)} & \piena{Elena} & \piena{Luca} & \piena{Sara} & \piena{Marco} & \texttt{testa 1, coda 1} \\
\cline{2-5}
\end{tabular}
\end{center}
\begin{itemize}\itemsep=2pt
    \item Anna esce dalla testa e Marco entra dal lato opposto, in fondo (posizione 3): la 0 resta libera; se gli indici avanzassero e basta, lo resterebbe per sempre
    \item \textbf{coda circolare}: dopo l'ultima posizione si torna alla 0 con il resto della divisione, \texttt{(coda + 1) \% elementi.length}: Elena va in 0 ``girando l'angolo'', ed è comunque l'ultima della fila
    \item quando i due indici coincidono (\texttt{testa == coda}) la coda può essere vuota o piena (ultima riga): le distingue \texttt{quanti}
\end{itemize}
\end{columns}
\end{frame}

% ========================================================================
\begin{frame}[fragile]{La coda con un array: lo stato e \texttt{enqueue} -- \texttt{CodaArray.java}}
\begin{columns}[T]
\column{0.46\textwidth}
\begin{lstlisting}[basicstyle=\ttfamily\scriptsize]
public class CodaArray<T>
        implements Coda<T> {

    private final Object[] elementi;
    private int testa;   // il primo
    private int coda;    // il prossimo
    private int quanti;

    public CodaArray(int capienza) {
        elementi = new Object[capienza];
        testa = 0;
        coda = 0;
        quanti = 0;
    }
\end{lstlisting}
\footnotesize
\begin{itemize}\itemsep=3pt
    \item \textbf{lo stato}: l'array, i due indici e il contatore. La regola che vale sempre: gli elementi, nell'ordine di arrivo, stanno in \texttt{testa}, \texttt{testa + 1}, \ldots{} fino a \texttt{coda - 1}, \textbf{girando l'angolo} se serve; sono \texttt{quanti}
    \item nel costruttore i tre \texttt{int} partono da 0: la coda è vuota, e il primo \texttt{enqueue} scrive in 0
    \item vuota e piena si leggono da \texttt{quanti}: in entrambi i casi i due indici coincidono
\end{itemize}

\column{0.54\textwidth}
\begin{lstlisting}[basicstyle=\ttfamily\scriptsize]
@Override
public boolean isEmpty() {
    return quanti == 0;
}

@Override
public boolean isFull() {
    return quanti == elementi.length;
}

@Override
public void enqueue(T x) {
    if (isFull()) {
        throw new IllegalStateException("piena");
    }
    elementi[coda] = x;
    coda = (coda + 1) % elementi.length;
    quanti++;
}
\end{lstlisting}
\footnotesize
\begin{itemize}\itemsep=3pt
    \item \texttt{enqueue} scrive dove indica \texttt{coda} e lo fa avanzare di uno, modulo la lunghezza: da \texttt{elementi.length - 1} si torna a \texttt{0}
\end{itemize}
\end{columns}
\end{frame}

% ========================================================================
\begin{frame}[fragile]{La coda con un array: \texttt{dequeue} -- \texttt{CodaArray.java}}
\begin{columns}[T]
\column{0.54\textwidth}
\begin{lstlisting}[basicstyle=\ttfamily\scriptsize]
@Override
public T dequeue() {
    if (isEmpty()) {
        throw new IllegalStateException("vuota");
    }
    T x = (T) elementi[testa];
    testa = (testa + 1) % elementi.length;
    quanti--;
    return x;
}
\end{lstlisting}
\footnotesize
\begin{itemize}\itemsep=3pt
    \item \texttt{dequeue} è simmetrico a \texttt{enqueue}: legge il primo, poi fa avanzare \texttt{testa}, sempre modulo la lunghezza; l'elemento resta fisicamente nell'array, ma fuori dalle posizioni usate: sarà sovrascritto da un \texttt{enqueue}
    \item il valore da restituire va salvato in \texttt{x} \textbf{prima} di spostare \texttt{testa}: dopo, \texttt{elementi[testa]} sarebbe il secondo
    \item lo stesso cast \texttt{(T)} della pila, con lo stesso avviso
\end{itemize}

\column{0.46\textwidth}
\begin{lstlisting}[basicstyle=\ttfamily\scriptsize]
Coda<String> c = new CodaArray<>(4);
c.enqueue("Anna");
c.enqueue("Luca");
c.enqueue("Sara");
System.out.println(c.dequeue());
c.enqueue("Marco");   // in posizione 3
c.enqueue("Elena");   // in posizione 0
System.out.println(c.isFull());
while (!c.isEmpty()) {
    System.out.println(c.dequeue());
}
\end{lstlisting}
\begin{lstlisting}[style=shell,basicstyle=\ttfamily\scriptsize\color{white}]
Anna
true
Luca
Sara
Marco
Elena
\end{lstlisting}
\footnotesize
\begin{itemize}\itemsep=3pt
    \item \texttt{ProvaCodaArray}: la sequenza della slide precedente; Elena sta in posizione 0 ma esce per ultima: conta l'ordine di arrivo
    \item il \texttt{while} con \texttt{isEmpty} è il modo giusto di svuotare: nessun \texttt{dequeue} senza aver chiesto
\end{itemize}
\end{columns}
\end{frame}

% ========================================================================
\begin{frame}[fragile]{Il nodo e la lista concatenata -- \texttt{Nodo.java}, \texttt{ListaConcatenata.java}}
\begin{columns}[T]
\column{0.5\textwidth}
\begin{lstlisting}[basicstyle=\ttfamily\scriptsize]
class Nodo<T> {

    T dato;
    Nodo<T> prossimo;

    Nodo(T dato) {
        this.dato = dato;
        prossimo = null;
    }
}
\end{lstlisting}
\begin{lstlisting}[basicstyle=\ttfamily\scriptsize]
public class ListaConcatenata<T> {

    private Nodo<T> testa;   // null se vuota

    public ListaConcatenata() {
        testa = null;
    }
\end{lstlisting}
\footnotesize
\begin{itemize}
    \item lista vuota: \texttt{testa == null}, ed è ciò che risponde \texttt{isEmpty()}. Una classe generica senza array né cast
    \item \textbf{niente capienza}: un nodo si crea per ogni elemento inserito, e quello di un elemento tolto diventa orfano (lezione~5)
\end{itemize}

\column{0.5\textwidth}
\begin{block}{Lista concatenata}
\small Una \emph{lista concatenata} (\emph{linked list}) è una sequenza di \emph{nodi}, ognuno con un dato e il riferimento al \emph{prossimo}; l'ultimo ha \texttt{prossimo} uguale a \texttt{null}. La lista conosce solo il primo nodo, la \emph{testa}.
\end{block}
\MermaidAdjustBoxOpts{max width=\linewidth,max height=0.2\textheight,center}
\begin{mermaid}
---
config:
  themeVariables:
    fontSize: 26px
  flowchart:
    nodeSpacing: 14
    rankSpacing: 30
    wrappingWidth: 600
---
flowchart LR
  t["<b>testa</b><br/><span style='color:#9ca3af'><i>riferimento</i></span>"]
  A["<b>Nodo</b><br/>dato: Anna<br/>prossimo<br/><span style='color:#9ca3af'><i>riferimento</i></span>"]
  B["<b>Nodo</b><br/>dato: Luca<br/>prossimo<br/><span style='color:#9ca3af'><i>riferimento</i></span>"]
  C["<b>Nodo</b><br/>dato: Sara<br/>prossimo<br/>null"]
  t --> A --> B --> C
  style t fill:#fef3c7,stroke:#b45309
\end{mermaid}
\vspace{-0.35cm}
\MermaidAdjustBoxOpts{max width=0.95\linewidth,max height=0.62\textheight,center}
\footnotesize
\begin{itemize}\itemsep=3pt
    \item \texttt{Nodo<T>} serve solo a \texttt{ListaConcatenata}: è dichiarata senza \texttt{public}, quindi la vedono solo le classi dello stesso package (lezione~9), e le sue variabili sono senza modificatore, così la lista le legge e le scrive direttamente (\texttt{q.prossimo = testa}), senza \texttt{get} e \texttt{set}; i client della lista non vedono mai un nodo: è dunque una \textbf{classe di servizio}
\end{itemize}
\end{columns}
\end{frame}


% ========================================================================
\begin{frame}[fragile]{Inserimento in testa e in coda -- \texttt{ListaConcatenata.java}}
\begin{columns}[T]
\column{0.46\textwidth}
\begin{lstlisting}[basicstyle=\ttfamily\scriptsize]
/** Inserisce x come primo elemento. */
public void addFirst(T x) {
    Nodo<T> q = new Nodo<>(x);
    q.prossimo = testa;
    testa = q;
}

/** Inserisce x come ultimo elemento. */
public void addLast(T x) {
    Nodo<T> q = new Nodo<>(x);
    if (isEmpty()) {
        testa = q;
    } else {
        Nodo<T> temp = testa;
        while (temp.prossimo != null) {
            temp = temp.prossimo;
        }
        temp.prossimo = q;
    }
}
\end{lstlisting}

\column{0.54\textwidth}
\footnotesize
\begin{lstlisting}[basicstyle=\ttfamily\scriptsize]
// nomi era [Anna, Luca]
nomi.addFirst("Marco");
\end{lstlisting}
\MermaidAdjustBoxOpts{max width=\linewidth,max height=0.3\textheight,center}
\begin{mermaid}
---
config:
  themeVariables:
    fontSize: 22px
  flowchart:
    nodeSpacing: 18
    rankSpacing: 36
    wrappingWidth: 600
---
flowchart LR
  t["<b>testa</b><br/><span style='color:#9ca3af'><i>riferimento</i></span>"]
  q["<b>q</b><br/><span style='color:#9ca3af'><i>riferimento</i></span>"]
  M["<b>Nodo</b><br/>dato: Marco<br/>prossimo"]
  A["<b>Nodo</b><br/>dato: Anna<br/>prossimo"]
  B["<b>Nodo</b><br/>dato: Luca<br/>prossimo: null"]
  q --> M
  M -- "1" --> A
  t -- "2" --> M
  t -. "prima" .-> A
  A --> B
  style t fill:#fef3c7,stroke:#b45309
  style M fill:#dcfce7,stroke:#15803d
\end{mermaid}
\MermaidAdjustBoxOpts{max width=0.95\linewidth,max height=0.62\textheight,center}
\begin{itemize}\itemsep=3pt
    \item \textbf{in testa}: 1) \texttt{q.prossimo = testa}, \texttt{q} punta al vecchio primo; 2) \texttt{testa = q}. L'\textbf{ordine conta}: con prima \texttt{testa = q} il resto della lista non sarebbe più raggiungibile. Su una lista vuota funziona lo stesso: \texttt{q.prossimo} resta \texttt{null}. Costa \textbf{un passo}
    \item \textbf{in coda}: \texttt{temp} scorre fino all'ultimo nodo, quello con \texttt{prossimo} uguale a \texttt{null}, e lo aggancia a \texttt{q}: costa quanto la lista è lunga; la lista vuota è un caso a parte: \texttt{q} diventa la testa
\end{itemize}
\end{columns}
\end{frame}

% ========================================================================
\begin{frame}[fragile]{Eliminazione in testa, interna e in coda -- \texttt{ListaConcatenata.java}}
\begin{columns}[T]
\column{0.56\textwidth}
\begin{lstlisting}[basicstyle=\ttfamily\scriptsize]
/** Toglie la prima x; vero se c'era. */
public boolean remove(T x) {
    if (isEmpty()) {
        return false;
    }
    if (testa.dato.equals(x)) {
        testa = testa.prossimo;   // era in testa
        return true;
    }
    Nodo<T> temp = testa;         // nodo prima di x
    while (temp.prossimo != null
            && !temp.prossimo.dato.equals(x)) {
        temp = temp.prossimo;
    }
    if (temp.prossimo == null) {
        return false;             // x non c'e'
    }
    temp.prossimo = temp.prossimo.prossimo;
    return true;
}
\end{lstlisting}
\footnotesize
\begin{itemize}\itemsep=2pt
    \item \textbf{in testa}: basta spostare \texttt{testa} al secondo nodo; il primo diventa orfano. \texttt{removeFirst()}, nel file, fa questo e restituisce il dato tolto (pre: non vuota), come il \texttt{pop} della pila
\end{itemize}

\column{0.44\textwidth}
\footnotesize
\begin{lstlisting}[basicstyle=\ttfamily\scriptsize]
// nomi era [Anna, Luca, Sara]
nomi.remove("Luca");
\end{lstlisting}
\MermaidAdjustBoxOpts{max width=\linewidth,max height=0.34\textheight,center}
\begin{mermaid}
---
config:
  themeVariables:
    fontSize: 26px
  flowchart:
    nodeSpacing: 14
    rankSpacing: 50
    wrappingWidth: 600
---
flowchart LR
  tm["<b>temp</b>"]
  A["dato: Anna<br/>prossimo"]
  B["dato: Luca<br/>prossimo"]
  C["dato: Sara<br/>prossimo: null"]
  tm --> A
  A -. "prima" .-> B
  B --> C
  A -- "dopo" --> C
  style B fill:#fee2e2,stroke:#b91c1c
\end{mermaid}
\MermaidAdjustBoxOpts{max width=0.95\linewidth,max height=0.62\textheight,center}
\begin{itemize}\itemsep=2pt
    \item \textbf{interna}: per scavalcare un nodo serve quello \textbf{prima}: \texttt{temp} si ferma sul nodo il cui prossimo contiene \texttt{x} (qui Anna), o sull'ultimo se \texttt{x} non c'è (l'\texttt{\&\&} non valuta \texttt{equals} quando \texttt{temp.prossimo} è \texttt{null}); la freccia ``dopo'' è \texttt{temp.prossimo = temp.prossimo.prossimo}
    \item \textbf{in coda}: se \texttt{x} è nell'ultimo nodo, \texttt{temp.prossimo.prossimo} è \texttt{null} e \texttt{temp} diventa l'ultimo: lo stesso codice
\end{itemize}
\end{columns}
\end{frame}

% ========================================================================
\begin{frame}[fragile]{La pila con una lista concatenata -- \texttt{PilaLista.java}}
\begin{columns}[T]
\column{0.54\textwidth}
\begin{lstlisting}[basicstyle=\ttfamily\scriptsize]
public class PilaLista<T>
        implements Pila<T> {

    private Nodo<T> testa;   // null se vuota
\end{lstlisting}
\begin{lstlisting}[basicstyle=\ttfamily\scriptsize]
@Override
public void push(T x) {
    Nodo<T> q = new Nodo<>(x);
    q.prossimo = testa;
    testa = q;
}

@Override
public T pop() {
    if (isEmpty()) {
        throw new IllegalStateException("vuota");
    }
    T x = testa.dato;
    testa = testa.prossimo;
    return x;
}
\end{lstlisting}

\column{0.46\textwidth}
\begin{lstlisting}[basicstyle=\ttfamily\scriptsize]
@Override
public boolean isEmpty() {
    return testa == null;
}

@Override
public boolean isFull() {
    return false;        // mai piena
}
\end{lstlisting}
{\footnotesize Il costruttore (\texttt{testa = null}) non è mostrato.}
\footnotesize
\begin{itemize}\itemsep=3pt
    \item \textbf{la cima della pila è la testa della lista}: \texttt{push} inserisce in testa, \texttt{pop} elimina in testa: un passo ciascuno, e di nuovo esce l'ultimo entrato, perché il nuovo nodo va sempre davanti ai vecchi
    \item \textbf{stessa interfaccia, altro corpo}: in \texttt{ProvaPile}, \texttt{Pila<String> q = new PilaLista<>()} con le stesse chiamate dà lo stesso output, tranne \texttt{isFull}: \texttt{false}, non c'è una capienza
\end{itemize}
\end{columns}
\end{frame}

% ========================================================================
\section{Esercitazione}

\begin{frame}[fragile]{Commentiamo insieme (20 min) -- il problema}
\begin{columns}[T]
\column{0.56\textwidth}
\footnotesize
\textbf{Il problema.} In un'espressione le parentesi tonde, quadre e graffe devono essere \textbf{bilanciate}: ogni parentesi chiusa deve chiudere l'ultima ancora aperta, ed essere dello stesso tipo; alla fine non deve restarne nessuna aperta; gli altri caratteri non contano. Il modulo \texttt{Parentesi} è già scritto: lo leggiamo riga per riga insieme a \texttt{ProvaParentesi}.

\vspace{0.1cm}
\textbf{Le regole}
\begin{itemize}\itemsep=2pt
    \item \texttt{isBilanciata(s)} restituisce \texttt{true} per \texttt{(a + b) * [c - \{d\}]} e per la stringa vuota
    \item \texttt{false} se una parentesi chiusa arriva quando non ce n'è nessuna aperta: \texttt{(a + b))}
    \item \texttt{false} se la chiusa non è del tipo dell'ultima aperta: \texttt{[(a + b])}
    \item \texttt{false} se alla fine qualcuna resta aperta: \texttt{((a + b)}
    \item lo strumento: una \textbf{pila} di caratteri, usata solo attraverso l'interfaccia \texttt{Pila<Character>}
\end{itemize}

\column{0.44\textwidth}
\footnotesize
\textbf{Perché una pila}: in \texttt{[(a + b])} la \texttt{]} deve chiudere l'\textbf{ultima} aperta, cioè \texttt{(}: ciò che è aperto per ultimo va chiuso per primo, LIFO
\begin{center}
\scriptsize\renewcommand{\arraystretch}{1.25}\setlength{\tabcolsep}{4pt}
\begin{tabular}{|c|l|l|}
\hline
\textbf{letto} & \textbf{azione} & \textbf{pila, cima a destra} \\
\hline
\texttt{[} & \texttt{push} & \texttt{[} \\
\hline
\texttt{(} & \texttt{push} & \texttt{[ (} \\
\hline
\texttt{a + b} & niente & \texttt{[ (} \\
\hline
\texttt{]} & \texttt{pop} dà \texttt{(} & \texttt{[} \\
\hline
 & \texttt{(} e \texttt{]} non corrispondono & \texttt{false} \\
\hline
\end{tabular}
\end{center}
\end{columns}
\end{frame}

% ========================================================================
\begin{frame}[fragile]{Commentiamo insieme -- \texttt{Parentesi.java}}
\begin{columns}[T]
\column{0.6\textwidth}
\begin{lstlisting}[basicstyle=\ttfamily\scriptsize]
/** Vero se le parentesi di s sono bilanciate. */
public static boolean isBilanciata(String s) {
    Pila<Character> pila = new PilaLista<>();
    for (int i = 0; i < s.length(); i++) {
        char c = s.charAt(i);
        if (c == '(' || c == '[' || c == '{') {
            pila.push(c);
        } else if (c == ')' || c == ']' || c == '}') {
            if (pila.isEmpty()
                    || !corrispondono(pila.pop(), c)) {
                return false;
            }
        }
    }
    return pila.isEmpty();   // nessuna aperta
}

private static boolean corrispondono(char a, char c) {
    return (a == '(' && c == ')')
            || (a == '[' && c == ']')
            || (a == '{' && c == '}');
}
\end{lstlisting}

\column{0.4\textwidth}
\footnotesize
\begin{itemize}\itemsep=4pt
    \item \texttt{Pila<Character>}, non \texttt{Pila<char>}: l'argomento di tipo generico è un tipo riferimento (lezione~11); \texttt{push(c)} inscatola il \texttt{char}, e \texttt{pila.pop()} passato a un parametro \texttt{char} viene tirato fuori (autoboxing e unboxing, lezione~4)
    \item \texttt{pila} è del tipo \textbf{interfaccia}: usa solo \texttt{push}, \texttt{pop}, \texttt{isEmpty}. Con \texttt{new PilaArray<>(s.length())} non cambierebbe nient'altro: più di \texttt{s.length()} parentesi non ci sono
    \item chiusa con la pila vuota, o diversa dall'ultima aperta: \texttt{false}. L'\texttt{||} non valuta il secondo operando se il primo è vero (lezione~6): \texttt{pop} solo su una pila non vuota
    \item alla fine, bilanciata se nessuna è rimasta aperta; \texttt{corrispondono} è \texttt{private} (lezione~4)
\end{itemize}
\end{columns}
\end{frame}

% ========================================================================
\begin{frame}[fragile]{Commentiamo insieme -- \texttt{ProvaParentesi.java}}
\begin{columns}[T]
\column{0.58\textwidth}
\begin{lstlisting}[basicstyle=\ttfamily\scriptsize]
String[] prove = { "(a + b) * [c - {d}]", "((a + b)",
                   "(a + b))", "[(a + b])", "" };
for (String s : prove) {
    System.out.println("\"" + s + "\" -> "
            + Parentesi.isBilanciata(s));
}
\end{lstlisting}
\begin{lstlisting}[style=shell,basicstyle=\ttfamily\scriptsize\color{white}]
"(a + b) * [c - {d}]" -> true
"((a + b)" -> false
"(a + b))" -> false
"[(a + b])" -> false
"" -> true
\end{lstlisting}
\footnotesize
\begin{itemize}\itemsep=3pt
    \item \texttt{"\textbackslash""} è il carattere \texttt{"} dentro una stringa (lezione~2)
    \item per la stringa vuota il \texttt{for} non parte e la pila è vuota: \texttt{true}
\end{itemize}

\column{0.42\textwidth}
\footnotesize
\textbf{Rispondiamo insieme:}
\begin{enumerate}\itemsep=3pt
    \item seguiamo \texttt{((a + b)}: che cosa c'è nella pila alla fine del \texttt{for}? Quale delle tre regole scatta?
    \item e \texttt{(a + b))}: a quale carattere si esce, e da quale \texttt{return}?
    \item se si togliesse \texttt{pila.isEmpty() ||} dalla condizione, che cosa succederebbe con \texttt{(a + b))}?
    \item perché \texttt{corrispondono} ha due parametri di tipo \texttt{char} e non \texttt{Character}? Che cosa succede nella chiamata \texttt{corrispondono(pila.pop(), c)}?
    \item con una coda al posto della pila, \texttt{[(a + b])} verrebbe accettata? Perché?
\end{enumerate}
\end{columns}
\end{frame}

% ========================================================================
\begin{frame}[fragile]{Insieme, passo passo (30 min) -- il problema}
\begin{columns}[T]
\column{0.58\textwidth}
\footnotesize
\textbf{Il problema.} Vogliamo una seconda realizzazione dell'ADT coda, \texttt{CodaLista<T>}, \textbf{senza capienza}: una lista concatenata da cui si toglie in testa (\texttt{dequeue}) e si aggiunge in fondo (\texttt{enqueue}). Il programma di prova \texttt{ProvaCode} è già scritto: simula la fila a uno sportello, prima con \texttt{CodaArray} e poi con \texttt{CodaLista}, e deve stampare le stesse righe, tranne una; poi usa \texttt{toString} su una coda di numeri.

\vspace{0.1cm}
\textbf{Le regole}
\begin{itemize}\itemsep=2pt
    \item \texttt{CodaLista<T>} implementa \texttt{Coda<T>}: \texttt{enqueue}, \texttt{dequeue}, \texttt{isEmpty}, \texttt{isFull}, con le pre e postcondizioni dell'interfaccia
    \item i nodi sono quelli di \texttt{Nodo<T>}; la classe non deve scorrere la lista per accodare: \texttt{enqueue} e \texttt{dequeue} costano un passo ciascuno
    \item \texttt{isFull} non è mai vera
    \item una precondizione violata solleva \texttt{IllegalStateException}
    \item \texttt{toString}: gli elementi dal primo all'ultimo, \texttt{[Luca, Sara]}
\end{itemize}

\column{0.42\textwidth}
\MermaidAdjustBoxOpts{max width=\linewidth,max height=0.55\textheight,center}
\begin{mermaid}
---
config:
  themeVariables:
    fontSize: 20px
  flowchart:
    nodeSpacing: 14
    rankSpacing: 30
    wrappingWidth: 600
---
flowchart TB
  P["<b>ProvaCode.java</b><br/>già scritto"]
  I["<b>Coda&lt;T&gt;</b><br/>l'interfaccia dell'ADT"]
  A["<b>CodaArray.java</b><br/>già scritta"]
  L["<b>CodaLista.java</b><br/>da scrivere"]
  N["<b>Nodo.java</b><br/>già scritto"]
  P -- "usa" --> I
  A -- "implementa" --> I
  L -- "implementa" --> I
  L -- "usa" --> N
  style P fill:#dcfce7,stroke:#15803d
  style A fill:#dcfce7,stroke:#15803d
  style N fill:#dcfce7,stroke:#15803d
  style I fill:#dbeafe,stroke:#1d4ed8
  style L fill:#fef3c7,stroke:#b45309
\end{mermaid}
\MermaidAdjustBoxOpts{max width=0.95\linewidth,max height=0.62\textheight,center}
\end{columns}
\end{frame}

% ========================================================================
\begin{frame}[fragile]{Insieme, passo passo -- \texttt{ProvaCode.java}}
\begin{columns}[T]
\column{0.6\textwidth}
\begin{lstlisting}[basicstyle=\ttfamily\scriptsize]
/** Chi arriva prima e' servito prima. */
public static void sportello(Coda<String> fila) {
    fila.enqueue("Anna");
    fila.enqueue("Luca");
    // fila era [Anna, Luca]
    fila.enqueue("Sara");
    System.out.println("piena: " + fila.isFull());
    System.out.println("servito " + fila.dequeue());
    fila.enqueue("Marco");
    while (!fila.isEmpty()) {
        System.out.println("servito " + fila.dequeue());
    }
}

public static void main(String[] args) {
    sportello(new CodaArray<>(3));
    sportello(new CodaLista<>());
    CodaLista<Integer> numeri = new CodaLista<>();
    numeri.enqueue(7);
    numeri.enqueue(4);
    numeri.enqueue(9);
    System.out.println(numeri.dequeue() + " " + numeri);
}
\end{lstlisting}

\column{0.4\textwidth}
{\footnotesize L'output atteso:}
\begin{lstlisting}[style=shell,basicstyle=\ttfamily\scriptsize\color{white}]
piena: true
servito Anna
servito Luca
servito Sara
servito Marco
piena: false
servito Anna
servito Luca
servito Sara
servito Marco
7 [4, 9]
\end{lstlisting}
\footnotesize
\begin{itemize}\itemsep=3pt
    \item \texttt{sportello} ha un parametro di tipo \textbf{interfaccia}: riceve l'una o l'altra coda e non se ne accorge; cambia solo la riga di \texttt{isFull}
    \item l'ultima riga usa \texttt{toString}, chiamato dalla concatenazione (lezione~8), dopo che il 7 è uscito
\end{itemize}
\end{columns}
\end{frame}

% ========================================================================
\begin{frame}[fragile]{Insieme, passo passo -- le decisioni}
\begin{columns}[T]
\column{0.62\textwidth}
\footnotesize
\textbf{Prima di scrivere il codice, decidiamo insieme:}
\begin{enumerate}\itemsep=3pt
    \item \textbf{lo stato?} \textrightarrow{} \textbf{due riferimenti}: \texttt{testa}, il primo nodo (quello che \texttt{dequeue} toglie) e \texttt{coda}, l'ultimo (dopo cui \texttt{enqueue} aggiunge); con la sola \texttt{testa}, \texttt{enqueue} dovrebbe scorrere tutta la lista come \texttt{addLast}. Coda vuota: entrambi \texttt{null}
    \item \textbf{\texttt{enqueue} su una coda vuota?} \textrightarrow{} non c'è un ultimo a cui agganciare \texttt{q}: \texttt{q} diventa la \texttt{testa}; altrimenti \texttt{coda.prossimo = q}. In entrambi i casi, poi, \texttt{coda = q}
    \item \textbf{\texttt{dequeue} dell'unico elemento?} \textrightarrow{} \texttt{testa} diventa \texttt{null}; anche \texttt{coda} va rimessa a \texttt{null}, per rispettare la regola dello stato (coda vuota: entrambi \texttt{null}): altrimenti continuerebbe a puntare a un nodo che non è più nella coda, e quel nodo non diventerebbe orfano
    \item \textbf{\texttt{isEmpty} e \texttt{isFull}?} \textrightarrow{} \texttt{testa == null}; sempre \texttt{false}
    \item \textbf{precondizioni?} \textrightarrow{} \texttt{dequeue} controlla \texttt{isEmpty} e solleva \texttt{IllegalStateException}; \texttt{enqueue} non ha nulla da controllare
    \item \textbf{\texttt{toString}?} \textrightarrow{} scorrere con \texttt{temp} da \texttt{testa} a \texttt{null}, come in \texttt{ListaConcatenata}
\end{enumerate}

\column{0.38\textwidth}
\MermaidAdjustBoxOpts{max width=\linewidth,max height=0.45\textheight,center}
\begin{mermaid}
classDiagram
  class CodaLista~T~ {
    -testa : Nodo~T~
    -coda : Nodo~T~
    +CodaLista()
    +enqueue(T)
    +dequeue() T
    +isEmpty() boolean
    +isFull() boolean
    +toString() String
  }
\end{mermaid}
\MermaidAdjustBoxOpts{max width=0.95\linewidth,max height=0.62\textheight,center}
\footnotesize
I metodi dell'interfaccia sono \texttt{public} e con \texttt{@Override} (lezione~9).
\end{columns}
\end{frame}

% ========================================================================
\begin{frame}[fragile]{Una possibile soluzione (1/2) -- lo stato e \texttt{enqueue}}
\begin{columns}[T]
\column{0.54\textwidth}
\begin{lstlisting}[basicstyle=\ttfamily\scriptsize]
public class CodaLista<T>
        implements Coda<T> {

    private Nodo<T> testa;   // il primo
    private Nodo<T> coda;    // l'ultimo

    public CodaLista() {
        testa = null;
        coda = null;
    }

    @Override
    public void enqueue(T x) {
        Nodo<T> q = new Nodo<>(x);
        if (isEmpty()) {
            testa = q;           // il primo
        } else {
            coda.prossimo = q;   // dopo l'ultimo
        }
        coda = q;            // ora e' l'ultimo
    }
\end{lstlisting}

\column{0.46\textwidth}
\begin{lstlisting}[basicstyle=\ttfamily\scriptsize]
// fila era [Anna, Luca]
fila.enqueue("Sara");
\end{lstlisting}
\MermaidAdjustBoxOpts{max width=\linewidth,max height=0.3\textheight,center}
\begin{mermaid}
---
config:
  themeVariables:
    fontSize: 22px
  flowchart:
    nodeSpacing: 18
    rankSpacing: 36
    wrappingWidth: 600
---
flowchart LR
  t["<b>testa</b>"]
  c["<b>coda</b>"]
  A["dato: Anna<br/>prossimo"]
  B["dato: Luca<br/>prossimo"]
  S["dato: Sara<br/>prossimo: null"]
  t --> A --> B
  B -- "1" --> S
  c -- "2" --> S
  c -. "prima" .-> B
  style S fill:#dcfce7,stroke:#15803d
\end{mermaid}
\MermaidAdjustBoxOpts{max width=0.95\linewidth,max height=0.62\textheight,center}
\begin{itemize}\itemsep=3pt
    \item 1) \texttt{coda.prossimo = q}, poi 2) \texttt{coda = q}; \texttt{q.prossimo} è già \texttt{null} dal costruttore di \texttt{Nodo}: il nuovo ultimo non punta a nulla
    \item il caso della coda vuota si distingue solo nell'aggancio; \texttt{coda = q} vale per tutti e due: non ripetiamolo nei due rami
    \item \texttt{enqueue} non controlla \texttt{isFull}: non può essere vera
\end{itemize}
\end{columns}
\end{frame}

% ========================================================================
\begin{frame}[fragile]{Una possibile soluzione (2/2) -- \texttt{dequeue}, i predicati, \texttt{toString}}
\begin{columns}[T]
\column{0.54\textwidth}
\begin{lstlisting}[basicstyle=\ttfamily\scriptsize]
@Override
public T dequeue() {
    if (isEmpty()) {
        throw new IllegalStateException("vuota");
    }
    T x = testa.dato;
    testa = testa.prossimo;
    if (testa == null) {
        coda = null;   // era l'unico
    }
    return x;
}

@Override
public boolean isEmpty() {
    return testa == null;
}

@Override
public boolean isFull() {
    return false;            // mai piena
}
\end{lstlisting}

\column{0.46\textwidth}
\begin{lstlisting}[basicstyle=\ttfamily\scriptsize]
@Override
public String toString() {
    String s = "[";
    Nodo<T> temp = testa;
    while (temp != null) {
        s += temp.dato;
        if (temp.prossimo != null) {
            s += ", ";
        }
        temp = temp.prossimo;
    }
    return s + "]";
}
\end{lstlisting}
\footnotesize
\begin{itemize}\itemsep=3pt
    \item \texttt{dequeue} è l'eliminazione in testa di \texttt{ListaConcatenata}, più l'\texttt{if} che azzera \texttt{coda} quando esce l'ultimo: lo stato torna quello della coda vuota, e \texttt{enqueue} la riconosce da \texttt{testa}
    \item \texttt{toString} scorre tutti i nodi: l'unico metodo che costa quanto la lunghezza della coda
    \item due realizzazioni, un'interfaccia: \texttt{sportello} non cambia
\end{itemize}
\end{columns}
\end{frame}

% ========================================================================
\section{Wrap-up}

\begin{frame}{Riepilogo: cosa abbiamo trattato}
\footnotesize
\begin{center}
\renewcommand{\arraystretch}{1.1}
\scriptsize
\begin{tabular}{|>{\raggedright\arraybackslash}p{6.4cm}|>{\raggedright\arraybackslash}p{7.4cm}|}
\hline
\textbf{Concetto (unità)} & \textbf{Dove compare oggi} \\
\hline
\texttt{Comparable} sulle nostre classi: l'ordine naturale (U23) & \texttt{Persona implements Comparable<Persona>}, \texttt{compareTo} per età \\
\hline
\texttt{Arrays.sort} su un array di oggetti; il vincolo \texttt{<T extends Comparable<T>>} (U23, U22) & \texttt{ProvaPersone}; \texttt{Utili.max(elenco)}; la \texttt{ClassCastException} senza \texttt{Comparable} \\
\hline
\texttt{Comparator}: un altro ordine, fuori dalla classe (U23) & \texttt{ComparaPerCognome}, \texttt{Arrays.sort(elenco, comparatore)} \\
\hline
Collezioni, contenitori, politiche di accesso; precondizioni e postcondizioni (U24) & la tabella delle politiche; i commenti di \texttt{Pila} e \texttt{Coda}; \texttt{IllegalStateException} \\
\hline
L'ADT pila, \emph{stack} (LIFO: \texttt{push}, \texttt{pop}) e l'ADT coda, \emph{queue} (FIFO: \texttt{enqueue}, \texttt{dequeue}); l'interfaccia dell'ADT (U24, U17) & \texttt{Pila<T>}, \texttt{Coda<T>}; \texttt{Parentesi} e \texttt{sportello} usano solo l'interfaccia \\
\hline
Pila e coda con un array; la coda circolare (U24) & \texttt{PilaArray}, \texttt{CodaArray} con \texttt{\% elementi.length} e \texttt{quanti} \\
\hline
Il nodo e la lista concatenata (\emph{linked list}): inserimento ed eliminazione in testa, in coda, interna (U24) & \texttt{Nodo<T>}, \texttt{ListaConcatenata<T>}; i diagrammi con \texttt{q} e \texttt{temp} \\
\hline
Pila e coda con una lista concatenata (U24) & \texttt{PilaLista}; \texttt{CodaLista} con \texttt{testa} e \texttt{coda} \\
\hline
\end{tabular}
\end{center}

\vspace{0.1cm}
\textbf{Prossima lezione (L13)}: i file di testo: stream di byte e di caratteri, gli stream standard, scrivere con \texttt{PrintWriter} e \texttt{FileWriter}, leggere con \texttt{Scanner} e \texttt{BufferedReader}, \texttt{close} e il try-with-resources (il \texttt{finally} della lezione~10); una rubrica su file caricata in una lista concatenata. Ripasso finale e simulazione della prova.
\end{frame}

% ========================================================================
\begin{frame}{Homework 12 (da svolgere entro la prossima lezione)}
\footnotesize\setlength{\parskip}{0pt}
Tre esercizi, ciascuno con il suo programma di prova, nella cartella \texttt{homework/}; dove servono, copiare \texttt{Nodo}, \texttt{Pila}, \texttt{PilaLista}, \texttt{Coda} e la \texttt{CodaLista} dell'esercitazione. Le precondizioni violate sollevano \texttt{IllegalStateException}.
\begin{enumerate}\itemsep=3pt
    \item \textbf{\texttt{Libro.java}} (con \texttt{ComparaPerAnno.java} e \texttt{ProvaLibri.java}): \texttt{Libro} è immutabile, con \texttt{titolo}, \texttt{autore} e \texttt{anno}, i tre \texttt{get} e \texttt{toString} nella forma \texttt{Effective Java (Bloch, 2018)}; implementa \texttt{Comparable<Libro>} con l'ordine naturale per \textbf{titolo}. \texttt{ComparaPerAnno} implementa \texttt{Comparator<Libro>}, per anno crescente. \texttt{ProvaLibri} crea un array di quattro libri, lo ordina con \texttt{Arrays.sort} e lo stampa un libro per riga; cerca con \texttt{Arrays.binarySearch} un \texttt{new Libro("Effective Java", "?", 0)} e stampa l'indice trovato (perché basta il titolo?); infine lo riordina con il comparatore e lo ristampa.
    \item \textbf{\texttt{ListaOrdinata.java}} (con \texttt{ProvaListaOrdinata.java}): \texttt{ListaOrdinata<T extends Comparable<T>>}, una lista concatenata di \texttt{Nodo<T>} i cui elementi restano sempre in ordine crescente (ripetizioni ammesse). \texttt{add(T x)} mette \texttt{x} al posto giusto con \textbf{un solo scorrimento}, fermandosi sull'ultimo nodo con dato non maggiore di \texttt{x} (in testa, interna o in coda: i casi di oggi); \texttt{remove(T x)} toglie la prima occorrenza e dice se c'era; \texttt{contains(T x)} si ferma appena incontra un dato maggiore di \texttt{x}; \texttt{isEmpty()}, \texttt{first()} (il più piccolo; pre: non vuota), \texttt{toString}. \texttt{ProvaListaOrdinata} inserisce i tempi \texttt{52, 48, 61, 48, 55} e stampa \texttt{[48, 48, 52, 55, 61]} e il migliore; rimuove un 48, il 61 e un valore assente, stampando ogni volta risultato e lista; prova \texttt{contains}; inserisce quattro nomi.
    \item \textbf{\texttt{Palindromi.java}} (con \texttt{ProvaPalindromi.java}): un modulo con \texttt{isPalindroma(String s)}, vero se le \textbf{lettere} di \texttt{s}, ignorando spazi, punteggiatura e maiuscole, si leggono uguali nei due versi. Ogni lettera, resa minuscola con \texttt{Character.toLowerCase}, va in una \texttt{Pila<Character>} e in una \texttt{Coda<Character>}; poi si svuotano insieme, confrontando ciò che esce dalla pila (al contrario) con ciò che esce dalla coda (in ordine): \texttt{false} alla prima differenza. \texttt{ProvaPalindromi} lo prova su \texttt{anna}, \texttt{I topi non avevano nipoti}, \texttt{Java} e \texttt{""}.
\end{enumerate}
\end{frame}

% ========================================================================
\begin{frame}
\begin{center}
{\Huge Domande?}\\[2em]
{\normalsize \href{mailto:mattia.fonisto@unina.it}{\texttt{mattia.fonisto@unina.it}}}
\end{center}
\end{frame}

\end{document}
